\documentclass[../../../../main.tex]{subfiles}
\begin{document}

\begin{flushright}
\footnotesize\textit{【自動文字起こし・要確認】}
\end{flushright}

{\bf [解]}
$0<\theta<90\quad\cdots①,\qquad a>0\ (a\in\mathbb R)\quad\cdots②$

(ii)から，$e(\theta)=\cos\theta^\circ+i\sin\theta^\circ$とおくと，
\begin{align*}
z_{n+1}-z_n=e(\theta)(z_n-z_{n-1})
\end{align*}

$z_0=0$，$z_1=a$とあわせて，くり返し用いて，
\begin{align*}
z_{n+1}-z_n=\{e(\theta)\}^na
\end{align*}

$\alpha_n=\dfrac{z_n}{\{e(\theta)\}^n}$とすると，
\begin{align*}
\alpha_{n+1}&=\frac1{e(\theta)}\alpha_n+\frac a{e(\theta)}\\
\alpha_{n+1}-\frac a{e(\theta)-1}&=\frac1{e(\theta)}\left\{\alpha_n-\frac a{e(\theta)-1}\right\}\quad(\because①から e(\theta)\ne1)
\end{align*}

$\alpha_0=0$とから，くり返し用いて，
\begin{align*}
\alpha_n&=\left(\frac1{e(\theta)}\right)^n\left(0-\frac a{e(\theta)-1}\right)+\frac a{e(\theta)-1}\\
z_n&=\frac a{e(\theta)-1}\left\{e(n\theta)-1\right\}
\end{align*}

（$\because$ド・モアブルの定理から$\{e(\theta)\}^n=e(n\theta)$）

これが$z_0=0$に一致する$n$（$n\ge1$）が存在する条件は，$a\ne0$，$e(\theta)-1\ne0$（$\because①②$）だから，
\begin{align*}
&\text{"}e(n\theta)=1\text{なる}n\ge1\text{が存在する"}\\
\iff&n\theta=360m\text{なる}n\in\mathbb N,\ m\in\mathbb Z\text{が存在する}\\
\iff&\theta=\frac mn\cdot360\text{なる}n\in\mathbb N,\ m\in\mathbb Z\text{が存在する}\\
\iff&\theta\text{が有理数}
\end{align*}

よって示された．$\blacksquare$

\newpage

\end{document}
